\section{Kに関する議論の経緯と現在の定義}

\subsection{本書の目的}

本書は、ECAにおける量 $K$ について、これまでに行った議論の経緯、そこから確定した構造、採用された定式化、および議論の過程で却下された考え方を整理することを目的とする。
ここで重要なのは、本書が新たな $K$ の定義を提案するためのものではないという点である。本書に記載する内容は、これまでのECAに関する議論によって確定した内容を、後から別の解釈へ変更されないように整理したものである。
特に、$K$ については、名称に含まれる「分散」という語から統計学的な分散を直接導入したり、一般的な数学的概念から距離や写像を先に設定したりするのではなく、ECAにおけるSolver集合の構造から出発して議論を進めてきた。
また、この「公理分散値」という命名が「統計学的な分散」を想起させてしまうことを鑑みて、準備が整い次第「Solver集合の公理配置指標」という命名に変更することも明記しておく。

したがって、本書では以下の点を明確に区別する。

\begin{itemize}
\item $K$ が最終的に何を記述する量なのか。
\item $K$ の定義に至るまでに、どのような集合構造を導入したのか。
\item 各数式がKの構造の中でどのような役割を持つのか。
\item どの考え方を採用し、どの考え方を却下したのか。
\item 却下した考え方が、なぜ現在のECAでは採用できないのか。
\end{itemize}

\subsection{Kについて最初に確定した基本方針}

Kについての議論で最も重要な確定事項は、Kを

\[
\boxed{
\text{Solver集合上における公理の配置・関係・分離を記述する量}
}
\]

として扱うことである。

ここでいう「分散」は、統計学におけるvarianceをそのまま意味するものではない。

Kが対象とするのは、Solver集合 $\mathcal{Z}$ の上に公理がどのように配置され、その配置同士がどのような関係を持ち、その結果としてどのような分離が生じるかという構造である。

したがって、Kの構造は、

\[
\boxed{
\text{配置}
\longrightarrow
\text{配置間の関係}
\longrightarrow
\text{分離}
\longrightarrow
\text{数値化}
}
\]

という順序で構成される。

この順序は単なる説明上の順序ではない。

まず公理の配置を定めなければ、配置同士を比較することはできない。また、配置を比較することによって関係が得られ、その関係の一つの特殊な形として分離が現れる。そして最後に、それらの構造を数値評価することによってKの値が得られる。

したがって、Kの数値を先に定めてからSolver集合上の構造を解釈するのではなく、

\[
\text{Solver集合上の公理配置}
\longrightarrow
\text{配置構造}
\longrightarrow
\text{数値評価}
\]

という方向でKを構成する必要がある。

\subsubsection{補足事項}
この「分散」という単語を「統計学に関連した何か」と断定する傾向はAIに共通するようだ。
しかし、このような解釈がAIから生まれるのは、分散という単語だけ強調して「「公理」が「どれくらい散らばっているか (分散)」を表す数値」と早合点して推測していることが根本的要因としてあるものと考えられる。
皮肉なことに、これはClaudeのみならずChatGPTにおいても同様の傾向が確認できた (率直に申し上げて、これはナンセンスと言わざるを得ない)。
したがって、この「公理分散値」という命名が「統計学的な分散」を想起させてしまうことを鑑みて、準備が整い次第「Solver集合の公理配置指標」という命名に変更する可能性についてあえて明記しておくこととした (準備が整うまでは命名の変更はしないものとする)。

\subsection{Kが対象とするSolver構造}

ECAでは、各Solver $z\in\mathcal{Z}$ に対して公理構造 $\mathcal{A}_z$ が対応する。

\[
z\in\mathcal{Z}
\longmapsto
\mathcal{A}_z
\]

ここで $\mathcal{A}_z$ はSolver $z$ に含まれる公理要素によって構成される。

一方、Solverには推論規則および構成制約を表す $\mathcal{R}_z$ も存在する。

\[
\mathcal{R}_z :
\text{推論規則および構成制約}
\]

しかし、Kの定義において $\mathcal{R}_z$ は直接の対象としない。

Kが記述する対象は、

\[
\mathcal{A}_z :
\text{Solver }z\text{ における公理構造}
\]

すなわち、Solver集合上における公理の配置だからである。

もし $\mathcal{R}_z$ までKの直接的な対象に含めるなら、Kが記述する対象は公理の配置だけではなく、推論規則や構成制約を含むSolverの構成全体へ拡張される。

それは、現在のKに与えられている役割とは異なる。

したがって、

\[
\boxed{
K\text{の直接的対象}=\mathcal{A}_z
}
\]

とし、

\[
\boxed{
\mathcal{R}_z\text{はKの直接的対象としない}
}
\]

という区別を確定した。

これは単に記述を簡略化するための選択ではなく、Kの対象範囲を「Solver集合上の公理配置」に限定するために必要な区別である。

\subsection{公理のSolver集合上での配置}

公理 $a$ に対して、その公理を含むSolverの集合を

\[
\boxed{
\mathcal{Z}_a = \{z\in\mathcal{Z}\mid a\in\mathcal{A}_z\}
}
\]

と定義する。

この式はKの構造において非常に重要である。

$\mathcal{Z}_a$ は単に「公理 $a$ を含むSolverの個数」を表すものではない。

$\mathcal{Z}_a$ は、

\[
\boxed{
\text{公理 }a\text{ がSolver集合 }\mathcal{Z}\text{ のどこに配置されているか}
}
\]

を集合として表現する。

したがって、$\mathcal{Z}_a$ はKの数値そのものではなく、Kが数値化する対象となる基礎構造である。

この定義によって、

\[
a
\longmapsto
\mathcal{Z}_a
\]

という対応が得られる。

さらに、これを明示的な写像として表現すると、

\[
\boxed{
\Phi_{\mathcal{Z}}:\mathcal{A}\longrightarrow\mathcal{P}(\mathcal{Z})
}
\]

\[
\boxed{
\Phi_{\mathcal{Z}}(a) = \mathcal{Z}_a = \{z\in\mathcal{Z}\mid a\in\mathcal{A}_z\}
}
\]

となる。

ここで $\mathcal{A}$ は対象となる公理全体の集合であり、$\mathcal{P}(\mathcal{Z})$ はSolver集合の冪集合である。

$\Phi_{\mathcal{Z}}$ の役割は、公理をSolver集合上の配置へ対応させることである。

したがって、

\[
\boxed{
\Phi_{\mathcal{Z}}\neq K
}
\]

である。

$\Phi_{\mathcal{Z}}$ は配置を表現するための構造的な写像であり、Kはその配置構造を記述し、最終的に数値化する量である。

この区別は重要である。写像そのものをKとみなしてしまうと、Kが集合なのか写像なのか数値なのかという対象の混同が生じるためである。

\subsection{複数公理の同時配置}

単一の公理だけでなく、複数の公理が同時に含まれるSolverを扱う必要がある。

そこで、$A\subseteq\mathcal{A}$ に対して、

\[
\boxed{
\mathcal{Z}_A = \bigcap_{a\in A}\mathcal{Z}_a
}
\]

を定義する。

$\mathcal{Z}_a$ の定義から、

\[
\boxed{
\mathcal{Z}_A = \{z\in\mathcal{Z}\mid A\subseteq\mathcal{A}_z\}
}
\]

が得られる。

この式の役割は、$A$ に含まれるすべての公理を同時に含むSolverを取り出すことである。

したがって、

\[
\mathcal{Z}_A
\]

は複数公理の同時配置をSolver集合上で表す集合となる。

特に、$a,b\in\mathcal{A}$ に対して、

\[
\boxed{
\mathcal{Z}_{{a,b}} = \mathcal{Z}_a\cap\mathcal{Z}_b
}
\]

となる。

これは、公理 $a$ と公理 $b$ が同時に配置されているSolverを直接取り出す構造である。

したがって、複数公理の関係を外部から仮定する必要はなく、それぞれの公理のSolver集合上の配置から直接構成できる。

\subsection{配置から公理間の関係を得る}

公理 $a,b\in\mathcal{A}$ に対して、

\[
\mathcal{Z}_a
\qquad\text{および}\qquad
\mathcal{Z}_b
\]

を比較することで、公理間の配置関係を得ることができる。

まず、

\[
\mathcal{Z}_a\cap\mathcal{Z}_b
\]

は、両方の公理が共通して配置されているSolverを表す。

一方、

\[
\mathcal{Z}_a\setminus\mathcal{Z}_b
\]

は、公理 $a$ は含むが $b$ は含まないSolverの配置を表し、

\[
\mathcal{Z}_b\setminus\mathcal{Z}_a
\]

は、公理 $b$ は含むが $a$ は含まないSolverの配置を表す。

さらに、両者の配置の差をまとめる場合には対称差

\[
\boxed{
\mathcal{Z}_a\triangle\mathcal{Z}_b = (\mathcal{Z}_a\setminus\mathcal{Z}_b) \cup (\mathcal{Z}_b\setminus\mathcal{Z}_a)
}
\]

を用いることができる。

これらはいずれも、まず $\mathcal{Z}_a$ と $\mathcal{Z}_b$ という配置が存在し、それらを比較することによって得られる集合構造である。

したがって、

\[
\boxed{
\text{配置}
\longrightarrow
\text{配置の比較}
\longrightarrow
\text{関係}
}
\]

という構造になる。

さらに、

\[
\mathcal{Z}_a\subseteq\mathcal{Z}_b
\]

であれば、

\[
\forall z\in\mathcal{Z},
\quad
a\in\mathcal{A}_z
\Longrightarrow
b\in\mathcal{A}_z
\]

が成り立つ。

これは、公理 $a$ が配置されているSolverでは必ず公理 $b$ も配置されているという関係を意味する。

したがって、包含関係もまた、公理間に外部から導入された関係ではなく、Solver集合上の配置から自然に得られる関係である。

\subsection{分離は配置関係から生じる}

公理 $a,b$ の配置が完全に重ならない場合、

\[
\boxed{
\mathcal{Z}_a\cap\mathcal{Z}_b=\varnothing
}
\]

となる。

これは、公理 $a$ と公理 $b$ がSolver集合上で完全に分離されていることを表す。

重要なのは、分離を配置とは独立した第三の対象として最初から導入しているわけではないことである。

分離は、

\[
\mathcal{Z}_a
\qquad\text{と}\qquad
\mathcal{Z}_b
\]

という二つの配置を比較した結果として現れる。

したがって、

\[
\boxed{
\text{配置}
\longrightarrow
\text{配置の比較}
\longrightarrow
\text{関係}
\longrightarrow
\text{分離}
}
\]

という構造になる。

特に、

\[
\mathcal{Z}_a\cap\mathcal{Z}_b=\varnothing
\]

は、この関係構造の中で完全な分離が現れた特殊な場合として位置付けられる。

このため、Kの構造において「配置」「関係」「分離」は独立した三つの選択肢ではない。

\subsection{配置の大きさと配置構造の区別}

$\mathcal{Z}_a$ の要素数

\[
|\mathcal{Z}_a|
\]

は、公理 $a$ を含むSolverの数を表す。

有限なSolver集合においては、

\[
|\mathcal{Z}_a|
\]

を調べることで、公理 $a$ がどの程度の数のSolverに現れるかを知ることができる。

しかし、

\[
|\mathcal{Z}_a| = |\mathcal{Z}_b|
\]

であっても、

\[
\mathcal{Z}_a = \mathcal{Z}_b
\]

とは限らない。

例えば、

\[
\mathcal{Z}_a=\{z_1,z_2\},
\qquad
\mathcal{Z}_b=\{z_3,z_4\}
\]

であれば、

\[
|\mathcal{Z}_a| = |\mathcal{Z}_b| = 2
\]

である一方、

\[
\mathcal{Z}_a\cap\mathcal{Z}_b = \varnothing
\]

である。

したがって、

\[
\boxed{
\text{配置の大きさ}
\neq
\text{配置そのもの}
}
\]

である。

これはKの数値化において重要な区別である。

例えば、後述する

\[
K(a)=\frac{|\mathcal{Z}_a|}{|\mathcal{Z}|}
\]

という数値化を考えた場合、この値は$\mathcal{Z}_a$の大きさを数値化するが、$\mathcal{Z}_a$がSolver集合上のどこに配置されているかという構造全体を完全に表現するものではない。

したがって、Kの基礎構造としては、

\[
\mathcal{Z}_a
\]

および

\[
\{\mathcal{Z}_a\mid a\in\mathcal{A}\}
\]

を保持し、その上で必要な数値評価を行う必要がある。

\subsection{Kの基本的定義}

以上を踏まえると、Kは次のように定義される。

\[
\boxed{
K\text{ は、Solver集合上に配置された公理について、その配置を基礎として公理間の関係および分離の構造を記述し、それを数値化する量である。}
}
\]

この定義では、Kは単に公理がいくつのSolverに存在するかだけを表す量ではない。

まず、

\[
\mathcal{Z}_a = \{z\in\mathcal{Z}\mid a\in\mathcal{A}_z\}
\]

によって公理の配置を定める。

次に、

\[
\mathcal{Z}_a\cap\mathcal{Z}_b,
\qquad
\mathcal{Z}_a\setminus\mathcal{Z}_b,
\qquad
\mathcal{Z}_b\setminus\mathcal{Z}_a,
\qquad
\mathcal{Z}_a\triangle\mathcal{Z}_b
\]

などによって配置間の関係を得る。

その中で、

\[
\mathcal{Z}_a\cap\mathcal{Z}_b=\varnothing
\]

という特殊な場合として完全な分離が現れる。

そして、最終的にこの配置構造を数値評価することでKの値が得られる。

したがって、Kの構造全体は、

\[
\boxed{
\mathcal{Z}_a
\longrightarrow
\{\mathcal{Z}_a\mid a\in\mathcal{A}\}
\longrightarrow
\text{配置間の関係}
\longrightarrow
\text{分離}
\longrightarrow
K
}
\]

として整理できる。

\subsection{KとB,Iとの関係}

KはSolver集合上の公理配置を数値化する量であるが、Kだけで公理体系全体の構造を表現するものではない。

ECAでは、B、I、Kはそれぞれ異なる側面を扱う。

Bは公理数や公理選択空間に関する量であり、Solver側の公理構成の規模を扱う。

Iは公理体系間の構造整合性を扱う。

これに対してKは、

\[
\text{公理}
\longrightarrow
\text{Solver集合上の配置}
\]

という方向から、公理がSolver集合上にどのように配置されているかを扱う。

したがって、

\[
B,\ I,\ K
\]

は同一の情報を別々の方法で表現するものではなく、それぞれ異なる構造を数値化する。

この区別によって、K単独では表現しきれない情報についても、ECAの数値評価

\[
\boxed{
T_{\mathrm{overlay}} = \frac{1}{N}\sum_{i=1}^{N}B_iI_iK_i
}
\]

においてBおよびIと組み合わせて扱うことができる。

したがって、

\[
\boxed{
K\text{単独の数値から公理配置の全構造を完全に復元できる必要はない}
}
\]

ということも確定している。

Kの基礎となる配置構造を保持しながら、それを数値として評価し、必要に応じてBおよびIと組み合わせてECAの数値評価へ接続する。

\subsection{定義から数値化への接続}

Kの定義段階では、

\[
\mathcal{Z}_a
\]

および

\[
\{\mathcal{Z}_a\mid a\in\mathcal{A}\}
\]

によってSolver集合上の公理配置を構成する。

その後、この配置構造に対して数値評価を導入する。

したがって、Kの定義とKの数値化は区別される。

\[
\boxed{
\text{集合構造}
\longrightarrow
\text{配置関係}
\longrightarrow
\text{数値評価}
}
\]

という流れである。

この区別によって、集合構造そのものと、それを評価して得られる数値とを混同せずに扱うことができる。

\subsection{Kの数値化における基本的な場合}

有限で単純なSolver集合を考え、各Solverを同等に扱う場合には、

\[
\boxed{
K(a)=\frac{|\mathcal{Z}_a|}{|\mathcal{Z}|}
}
\]

という数値化を行うことができる。

この式において、

\[
|\mathcal{Z}_a|
\]

は公理 $a$ を含むSolverの数、

\[
|\mathcal{Z}|
\]

は対象となるSolver全体の数である。

したがって、この値は、公理 $a$ のSolver集合全体に対する配置の規模を正規化した値として解釈できる。

ただし、この式そのものをKの一般的な定義と同一視してはいけない。

Kの定義は、

\[
\boxed{
\text{Solver集合上の公理配置・配置間の関係・分離を記述し、それを数値化する量}
}
\]

であり、

\[
K(a)=\frac{|\mathcal{Z}_a|}{|\mathcal{Z}|}
\]

は、その構造を数値化する具体的な方法の一つである。

この区別は、有限かつ一様なSolver集合だけでなく、より一般のSolver集合を扱うためにも必要である。

\subsection{一般的な数値評価への拡張}

より一般の数値評価を考える場合には、まず配置構造 $\mathcal{Z}_a$ に対して評価を与える。

\[
\boxed{
\widetilde{K}(a)=W_K(\mathcal{Z}_a)
}
\]

ここで $W_K$ は、対象とするSolver集合の構造や導入する測度等に応じて、配置を評価するための量である。

さらに、その値を適切に正規化する場合には、

\[
\boxed{
K(a) = \frac{\widetilde{K}(a)}{\sup_{b\in\mathcal{F}_K}\widetilde{K}(b)}
}
\]

という形を用いることができる。

ここで重要なのは、$\widetilde{K}(a)$ や $K(a)$ が新たな配置構造を作るものではないことである。

基礎にあるのはあくまで

\[
\mathcal{Z}_a
\]

であり、

\[
W_K
\]

はその配置を数値評価するための方法である。

したがって、

\[
\boxed{
\mathcal{Z}_a
\neq
\widetilde{K}(a)
\neq
K(a)
}
\]

であり、それぞれ、

\[
\mathcal{Z}_a
\quad\text{：配置そのもの}
\]

\[
\widetilde{K}(a)
\quad\text{：配置に対する未正規化の数値評価}
\]

\[
K(a)
\quad\text{：正規化されたKの数値}
\]

という役割を持つ。

\subsection{議論の中で却下された考え方}

\subsubsection{統計学的な分散をそのままKとする考え方}

Kについて初期に問題となったのは、「分散」という名称から統計学的なvarianceを直接導入する考え方である。

例えば、何らかの量 $d$ に対して

\[
\operatorname{Var}(d)
\]

を計算し、それをKとするような構成である。

しかし、この方向は現在のECAでは採用できない。

理由は、Kが最初に扱う対象が統計データではなく、Solver集合上の公理配置だからである。

現在のECAでは、

\[
\mathcal{Z}_a = \{z\in\mathcal{Z}\mid a\in\mathcal{A}_z\}
\]

によって公理の配置を構成することが先にある。

その配置から、

\[
\mathcal{Z}_a\cap\mathcal{Z}_b,
\qquad
\mathcal{Z}_a\setminus\mathcal{Z}_b,
\qquad
\mathcal{Z}_a\triangle\mathcal{Z}_b
\]

などの関係を得て、そこから分離を記述し、最後に数値化する。

したがって、

\[
\text{統計学的variance}
\longrightarrow
K
\]

という順序ではなく、

\[
\text{公理配置}
\longrightarrow
\text{配置関係}
\longrightarrow
\text{分離}
\longrightarrow
\text{数値化}
\]

という順序が必要である。

したがって、Kの「分散」という名称を統計学的varianceの定義そのものとして扱うことは却下された。

\subsubsection{最初から公理間の距離を定義する考え方}

別の方向として、公理 $a,b$ の間に最初から距離

\[
d(a,b)
\]

を定義し、その距離からKを構成する考え方も検討された。

しかし、現在のECAでは、公理集合 $\mathcal{A}$ に対してあらかじめ距離構造が与えられているわけではない。

したがって、

\[
d:\mathcal{A}\times\mathcal{A}\longrightarrow\mathbb{R}
\]

のような距離を最初から仮定することは、現在のECAのSolver集合構造から自然に導かれたものではない。

現在の構造では、まず

\[
\mathcal{Z}_a
\qquad\text{および}\qquad
\mathcal{Z}_b
\]

を構成する。

その上で、

\[
\mathcal{Z}_a\cap\mathcal{Z}_b
\]

や

\[
\mathcal{Z}_a\triangle\mathcal{Z}_b
\]

などの配置関係が得られる。

必要であれば、これらの関係を後段で数値化することは可能である。

しかし、距離をKの出発点として導入することは、

\[
\text{Solver集合上の配置}
\]

から始めるという現在のECAの構造を逆転させてしまう。

したがって、

\[
\boxed{
\text{距離をKの出発点とすることは採用しない}
}
\]

とした。

\subsubsection{写像 $\Phi_{\mathcal{Z}}$ そのものをKとする考え方}

\[
\Phi_{\mathcal{Z}}:\mathcal{A}\longrightarrow\mathcal{P}(\mathcal{Z})
\]

は、公理をSolver集合上の配置へ対応させるための写像である。

しかし、

\[
\Phi_{\mathcal{Z}}(a)=\mathcal{Z}_a
\]

はあくまで配置を表現する構造であり、Kそのものではない。

Kは、この配置構造を基礎として関係や分離を記述し、最終的に数値化する量として定義されている。

したがって、

\[
\boxed{
\Phi_{\mathcal{Z}}\neq K
}
\]

である。

この区別を失うと、「Kが集合なのか、写像なのか、数値なのか」という異なるレベルの対象を混同することになる。

\subsubsection{$\mathcal{R}_z$までKに含める考え方}

各Solverには、

\[
\mathcal{A}_z
\]

だけでなく、

\[
\mathcal{R}_z
\]

も存在する。

しかし、Kの対象を$\mathcal{R}_z$まで拡張することは採用しない。

$\mathcal{R}_z$は推論規則および構成制約を表すため、それをKに含めるとKが記述する対象が、

\[
\text{公理配置}
\]

から

\[
\text{Solverの構成全体}
\]

へ拡張されるからである。

これはKの現在の役割と一致しない。

したがって、Kでは

\[
\boxed{
\mathcal{A}_z\text{を直接の基礎対象とし、}\mathcal{R}_z\text{は直接扱わない}
}
\]

という区別を維持する。

\subsubsection{Kを単なる配置の大きさとする考え方}

有限なSolver集合では、

\[
K(a)=\frac{|\mathcal{Z}_a|}{|\mathcal{Z}|}
\]

という数値化が可能である。

しかし、これをもって「Kとは公理を含むSolverの割合そのものである」と定義することはできない。

なぜなら、

\[
|\mathcal{Z}_a|

|\mathcal{Z}_b|
\]

であっても、

\[
\mathcal{Z}_a

\mathcal{Z}_b
\]

とは限らないからである。

配置の大きさだけでは、

\[
\mathcal{Z}_a\cap\mathcal{Z}_b,
\qquad
\mathcal{Z}_a\setminus\mathcal{Z}_b,
\qquad
\mathcal{Z}_a\triangle\mathcal{Z}_b
\]

などの構造を表現できない。

したがって、

\[
\frac{|\mathcal{Z}_a|}{|\mathcal{Z}|}
\]

はKの数値化方法の一つではあるが、Kの構造全体をそれだけで定義するものではない。

\subsubsection{K単独の数値から配置構造を完全に復元させる考え方}

Kは最終的に数値として扱われるが、一つの数値から元の配置構造を完全に復元できなければならない、という条件は現在のECAには存在しない。

例えば、

\[
K(a)=K(b)
\]

であっても、

\[
\mathcal{Z}_a=\mathcal{Z}_b
\]

とは限らない。

同じKの値を持つ二つの公理が、異なるSolver集合上に配置されている可能性があるためである。

Kの基礎には、

\[
\mathcal{Z}_a
\]

および

\[
\{\mathcal{Z}_a\mid a\in\mathcal{A}\}
\]

という集合構造があり、Kはその構造を数値評価する。

したがって、

\[
\boxed{
\text{数値Kから配置構造全体を完全に復元できる必要はない}
}
\]

とする。

また、このことは、KをBおよびIと組み合わせて

\[
T_{\mathrm{overlay}} = \frac{1}{N}\sum_{i=1}^{N}B_iI_iK_i
\]

として扱うこととも矛盾しない。

B、I、Kは同じ構造を異なる方法で表すものではなく、それぞれ異なる側面を数値化する量である。

\subsection{現在のKにおける「分散」の意味}

Kにおける「分散」は、統計学的な分散値を意味するものとして定義されていない。

Kが記述する「分散」は、

\[
\boxed{
\text{Solver集合上における公理の配置・関係・分離}
}
\]

として理解する。

例えば、公理 $a$ と公理 $b$ に対して、

\[
\mathcal{Z}_a\cap\mathcal{Z}_b
\]

が大きい場合と、

\[
\mathcal{Z}_a\cap\mathcal{Z}_b=\varnothing
\]

の場合では、Solver集合上の配置関係が異なる。

また、

\[
\mathcal{Z}_a\subseteq\mathcal{Z}_b
\]

の場合には、公理 $a$ の配置が公理 $b$ の配置に含まれるという構造が存在する。

このような配置の違いを記述することが、Kにおける「分散」の意味である。

したがって、

\[
\boxed{
K\text{は統計学的varianceを別名で呼んだものではない}
}
\]

という点を明確にする必要がある。

\subsection{現在確定しているKの構造}

これまでの議論によって、Kについて以下の構造が確定している。

\begin{enumerate}
  \item 各Solver $z\in\mathcal{Z}$ には公理構造 $\mathcal{A}_z$ が対応する
  \item Kの直接的対象は $\mathcal{A}_z$ であり、$\mathcal{R}_z$ は直接の対象としない
  \item 公理 $a$ のSolver集合上の配置を $\mathcal{Z}_a = \{z\in\mathcal{Z}\mid a\in\mathcal{A}_z\}$ として定義する
  \item $\Phi_{\mathcal{Z}}:\mathcal{A}\to\mathcal{P}(\mathcal{Z})$ は、公理をSolver集合上の配置へ対応させる写像であり、Kそのものではない
  \item 複数公理 $A\subseteq\mathcal{A}$ の同時配置を $\mathcal{Z}_A = \bigcap_{a\in A}\mathcal{Z}_a$ として表す
  \item $\mathcal{Z}_a$ と $\mathcal{Z}_b$ の比較から、共通配置、差、対称差、包含などの配置関係が得られる
  \item $\mathcal{Z}_a\cap\mathcal{Z}_b=\varnothing$ は配置関係から生じる完全分離を表す
  \item 配置の大きさ「$|\mathcal{Z}_a|$」と配置構造「$\mathcal{Z}_a$」は区別する
  \item 有限かつ単純な場合には $K(a)=\frac{|\mathcal{Z}_a|}{|\mathcal{Z}|}$ によって配置を数値化できる
  \item 一般の場合には、配置に対する評価 $\widetilde{K}(a)=W_K(\mathcal{Z}_a)$ および必要な正規化によってKを数値化できる
  \item KはSolver集合上の公理配置・配置間の関係・分離を数値化するECA固有の量であること
  \item Kは統計学的varianceそのものでもなければ、統計学に関する概念でもないこと、およびECAが統計学に関連する何かでもないこと
  \item K単独の数値が公理配置の全構造を完全に復元する必要はない
  \item KはBおよびIとは異なる構造を数値化する量であり、$T_{\mathrm{overlay}} = \frac{1}{N}\sum_{i=1}^{N}B_iI_iK_i$ において組み合わせて扱うことができる
\end{enumerate}

\subsection{Kに関する最終的な位置付け}

以上の議論をまとめると、Kは次の構造を持つ。

\[
\boxed{
z
\longmapsto
\mathcal{A}_z
\longmapsto
\mathcal{Z}_a
\longmapsto
\{\mathcal{Z}_a\mid a\in\mathcal{A}\}
\longmapsto
\text{配置間の関係}
\longmapsto
\text{分離}
\longmapsto
K
}
\]

ここで各段階の役割は異なる。

\[
\mathcal{A}_z
\]

は各Solverにおける公理構造である。

\[
\mathcal{Z}_a
\]

は、公理 $a$ がSolver集合上のどこに配置されているかを表す。

\[
\{\mathcal{Z}_a\mid a\in\mathcal{A}\}
\]

は、対象となる公理全体のSolver集合上の配置構造を表す。

それらを比較することによって、

\[
\mathcal{Z}_a\cap\mathcal{Z}_b,
\qquad
\mathcal{Z}_a\setminus\mathcal{Z}_b,
\qquad
\mathcal{Z}_a\triangle\mathcal{Z}_b,
\qquad
\mathcal{Z}_a\subseteq\mathcal{Z}_b
\]

などの配置関係が得られる。

さらに、

\[
\mathcal{Z}_a\cap\mathcal{Z}_b=\varnothing
\]

という特殊な関係として完全な分離が現れる。

そして、これらの集合構造を適切な方法によって数値評価することでKの値が得られる。

したがって、現在のECAにおけるKは、

\[
\boxed{
K=
\text{Solver集合上における公理の配置・関係・分離を基礎として、それを数値化する量}
}
\]

として確定している。

\subsection{Claudeによる再検討における注意事項}

本書の内容を今後の検討に利用する場合、以下の点を維持する必要がある。

\begin{enumerate}
  \item ECAにおいてはKを統計学的varianceとして定義、解釈しないこと。
  \item 公理間の距離を最初から仮定してKの定義を開始しないこと。
  \item $\Phi_{\mathcal{Z}}$をKそのものと解釈しないこと。
  \item $\mathcal{R}_z$をKの直接的対象へ拡張しないこと。
  \item $|\mathcal{Z}_a|$だけをKの定義そのものとみなさないこと。
  \item Kの単一の数値に公理配置の全構造を完全に符号化することを要求しないこと。
  \item 既に確定した「配置 $\rightarrow$ 関係 $\rightarrow$ 分離 $\rightarrow$ 数値化」という順序を逆転させないこと。
  \item Kについて新しい改善案を提示する場合には、まず現在の定義に実際の欠陥が存在するかを確認し、単なる別案と必須修正を区別すること。
\end{enumerate}

特に最後の点は重要である。

現在のKは、これまでの議論によって上記の構造として確定したものである。
したがって、既存の定義とは異なる数学的構成を思いついたとしても、それだけを理由に現在の定義を変更する必要はない。

新しい構成を検討する場合には、

\[
\boxed{
\text{現在のKに実際の矛盾・誤り・定義不足があるのか}
}
\]

をまず確認する必要がある。

単に「別の方法も考えられる」ということと、「現在の方法が誤っている」ということは同じではない。

本書の目的は、この区別を含めてKに関するこれまでの議論を保存し、今後の検討において現在のECAの構造と議論の経緯が失われないようにすることである。